\documentclass[11pt,a4paper]{article}

% ── Compilar con: tectonic formulas.tex  (motor XeTeX, UTF-8 nativo) ──
\usepackage{fontspec}
\usepackage{amsmath,amssymb}
\usepackage[margin=2.3cm]{geometry}
\usepackage{xcolor}
\usepackage{booktabs}
\usepackage{tabularx}
\usepackage{listings}
\usepackage{enumitem}
\usepackage{underscore}
\usepackage{titlesec}
\usepackage[hidelinks]{hyperref}

% ── Colores (mismos que docs/plans/algoritmos.tex) ──
\definecolor{cazul}{HTML}{2C5F8A}
\definecolor{cverde}{HTML}{2E8B57}
\definecolor{cnaranja}{HTML}{C1660A}
\definecolor{cgris}{HTML}{5B6770}
\definecolor{cfondo}{HTML}{EDF2F7}
\definecolor{ccodefondo}{HTML}{F6F8FA}

\setcounter{secnumdepth}{3}
\setcounter{tocdepth}{2}

\renewcommand{\contentsname}{Índice}
\renewcommand{\tablename}{Tabla}

\lstset{
  basicstyle=\ttfamily\footnotesize,
  backgroundcolor=\color{ccodefondo},
  frame=single,
  rulecolor=\color{cgris!40},
  breaklines=true,
  columns=fullflexible,
  keepspaces=true,
  showstringspaces=false,
  xleftmargin=6pt,
  framexleftmargin=4pt,
}
\newcommand{\code}[1]{\texttt{\small #1}}

% ── Caja de nota ──
\newenvironment{nota}{%
  \par\medskip\noindent
  \begin{tabular}{@{}p{0.4em}@{\hspace{8pt}}p{\dimexpr\linewidth-2em\relax}@{}}
  \textcolor{cnaranja}{\rule{3pt}{1.1em}} &
}{\end{tabular}\par\medskip}

\newcommand{\bigsep}{\par\vspace{3pt}}

\title{\vspace{-1.4cm}\textbf{Fórmulas de activación de X1}\\[6pt]
  \large O0 $\to$ OE \quad$\cdot$\quad primer Stop Loss ganador\\[4pt]
  \normalsize\color{cgris} Derivación algebraica y ejemplos numéricos por activo}
\author{Proyecto Alginvesting}
\date{25 de septiembre de 2026}

\begin{document}
\maketitle
\thispagestyle{empty}

\begin{center}
\begin{minipage}{0.86\linewidth}
\small\color{cgris}
Este documento deriva, solo con álgebra, los dos precios de activación más consultados
de \code{X1_trading.py}: a qué precio un soporte (O0) recibe su buy limit (pasa a OE), y
a qué precio una posición abierta (OA) recibe su primer Stop Loss (SL) — y en qué precio
queda ese SL. Cierra con una tabla numérica por cada uno de los 6 activos operados.
Contexto: revisión de sesión del 25-09-2026, ver \code{docs/context/decisiones.md} en
esa fecha para el porqué del diseño actual.
\end{minipage}
\end{center}

\vspace{8pt}
\tableofcontents
\newpage

% ═══════════════════════════════════════════════════════════════════════════
\section{Notación}
% ═══════════════════════════════════════════════════════════════════════════

\begin{table}[h]
\centering
\small
\begin{tabularx}{\linewidth}{@{}l X l@{}}
\toprule
\textbf{Símbolo} & \textbf{Significado} & \textbf{Dónde vive} \\
\midrule
$P_0$ & Precio actual (bid). & tick en vivo, MT5 \\
$P_i$ & Precio de referencia: el del soporte en el Caso~1, el de apertura de la OA en el Caso~2. & JSON de soportes / \code{orden.price\_open} \\
$A$ & Umbral en USD (gap de activación / ganancia mínima). & \code{config.py} — \code{A[valor]} \\
$B$ & Distancia en USD que el trailing mantiene bajo $P_0$. & \code{config.py} — \code{B[valor]} \\
$U$ & Unidades del activo por lote. & \code{config.py} — \code{UNITS[valor]} \\
$M$ & Multiplicador de lote configurado (\code{LOTAJES\_M}). & \code{config.py} — \code{LOTAJES\_M[valor]} \\
$m$ & Lote mínimo del bróker. & \code{config.py} — \code{MIN\_LOTAJES[valor]} \\
$L$ & Exposición en USD por unidad de precio (\emph{qué lotaje} depende del caso). & calculado \\
\bottomrule
\end{tabularx}
\caption{Símbolos usados en ambos casos.}
\end{table}

\begin{nota}
$M$ y $m$ no son lo mismo: $m$ es una propiedad fija del bróker/activo (el lote más
chico que se puede operar); $M$ es una decisión de trading (cuántos lotes mínimos se
opera por posición). El lotaje configurado es \code{LOTAJES[valor]} $= M \cdot m$.
\end{nota}

% ═══════════════════════════════════════════════════════════════════════════
\section{Caso 1 --- O0 se convierte en OE}
% ═══════════════════════════════════════════════════════════════════════════

\subsection{Regla en código}

\code{crear\_ordenes\_espera} (\code{X1\_trading.py:421}) recorre cada soporte $P_i$ de
\code{lista\_N} y se lo salta —no lo declara como buy limit— si no cumple el gap
mínimo:

\begin{lstlisting}
if (P0 - Pi) * MIN_LOTAJES[valor] * UNITS[valor] < a:
    continue
\end{lstlisting}

Es decir, se \emph{crea} la OE cuando la condición de arriba es falsa: cuando
$(P_0 - P_i)\, m\, U \geq A$.

\subsection{Derivación}

\begin{align}
(P_0 - P_i)\, m\, U &\geq A \\
P_0 - P_i &\geq \frac{A}{mU} \\
P_0 &\geq P_i + \frac{A}{mU}
\end{align}

\begin{center}
\colorbox{cfondo}{$\displaystyle P_{0}^{\,O0\to OE} \;=\; P_i + \frac{A}{mU}$}
\end{center}

\begin{nota}
$M$ \textbf{no aparece.} Antes de la sesión del 25-09-2026 la fórmula usaba
$L = M m U$ (el lotaje \emph{configurado}, no el mínimo) — subir $M$ relajaba el gap de
activación sin motivo, porque ese gap es una propiedad del activo (qué tan lejos del
soporte vale la pena declarar la orden), no del tamaño de la posición que se abriría
ahí. Ver \code{docs/context/decisiones.md}, entrada 2026-09-25.
\end{nota}

\subsection{Ejemplos por activo}

$P_i$ es ilustrativo (un nivel de precio razonable por activo, no una cotización en
vivo) salvo AMZN, que usa el precio real revisado en la sesión del 25-09-2026
(captura de MT5, posición abierta en 248,37).

\begin{table}[h]
\centering
\small
\begin{tabularx}{\linewidth}{@{}l r r r r r r r@{}}
\toprule
\textbf{Activo} & $m$ & $U$ & $A$ & $mU$ & $A/(mU)$ & $P_i$ & $P_0^{\,O0\to OE}$ \\
\midrule
BTCUSD & 0,01 & 1   & 3 & 0,01 & 300,00 & 65\,000,00 & \textbf{65\,300,00} \\
ETHUSD & 0,10 & 1   & 3 & 0,10 & 30,00  & 3\,000,00  & \textbf{3\,030,00}  \\
TSLA   & 0,01 & 100 & 3 & 1,00 & 3,00   & 250,00     & \textbf{253,00}     \\
GOOGL  & 0,01 & 100 & 3 & 1,00 & 3,00   & 170,00     & \textbf{173,00}     \\
NVDA   & 0,01 & 100 & 3 & 1,00 & 3,00   & 130,00     & \textbf{133,00}     \\
AMZN   & 0,01 & 100 & 3 & 1,00 & 3,00   & 248,37     & \textbf{251,37}     \\
\bottomrule
\end{tabularx}
\caption{Precio al que cada soporte $P_i$ recibe su buy limit (pasa de O0 a OE). Config.
vigente: $A=3$ para los 6 activos.}
\end{table}

Lectura: un soporte de BTCUSD en 65\,000 recién se declara OE cuando el precio sube a
65\,300 (gap de \$300, porque $U=1$ multiplica el movimiento crudo de precio 1:1); el
mismo gap para una acción ($U=100$) es de solo \$3 de precio, porque cada \$1 de
movimiento ya representa \$100 de exposición por lote mínimo.

% ═══════════════════════════════════════════════════════════════════════════
\section{Caso 2 --- Primera activación del Stop Loss}
% ═══════════════════════════════════════════════════════════════════════════

\subsection{Regla en código}

\code{trailing\_stop} (\code{X1\_trading.py:541-553}) opera sobre el lotaje
\textbf{real} de la orden ya abierta (\code{orden.volume}), no sobre el lotaje
configurado:

\begin{lstlisting}
lotajes_m_efectivo = orden.volume / MIN_LOTAJES[valor]
L = orden.volume * UNITS[valor]
a_efectivo = a * lotajes_m_efectivo
b_efectivo = b * lotajes_m_efectivo
sl_nuevo = P0 - b_efectivo / L
...
if sl == 0:
    if ganancia >= a_efectivo:      # ganancia = (P0 - Pi) * L
        cambios = cambiar_SL(orden, valor, sl_nuevo, silent=True)
\end{lstlisting}

\subsection{Derivación del umbral de activación}

Si la orden se abrió al lotaje configurado, $\text{orden.volume} = Mm$:

\begin{align}
\text{lotajes\_m\_efectivo} &= \frac{Mm}{m} = M \\
L &= Mm \cdot U \\
a_{\text{efectivo}} &= A \cdot M \\
\text{ganancia} &= (P_0 - P_i)\, L = (P_0 - P_i)\, M m U
\end{align}

Condición de activación ($\text{ganancia} \geq a_{\text{efectivo}}$):

\begin{align}
(P_0 - P_i)\, M m U &\geq A M \qquad \text{(dividiendo por } M>0 \text{)} \\
(P_0 - P_i)\, m U &\geq A \\
P_0 &\geq P_i + \frac{A}{mU}
\end{align}

\begin{center}
\colorbox{cfondo}{$\displaystyle P_{0}^{\,SL_1} \;=\; P_i + \frac{A}{mU}$}
\end{center}

\begin{nota}
$M$ se cancela: aparece igual en \code{ganancia} y en \code{a\_efectivo}, así que el
umbral de precio no depende del lotaje de la posición. El resultado es
\textbf{algebraicamente idéntico} al del Caso~1 — misma fórmula, con $P_i$ = precio de
apertura de la OA en vez de precio del soporte. Antes de la sesión del 25-09-2026 esto
no era exacto (el Caso~1 sí tenía $M$ en el denominador); ahora ambas reglas son
literalmente la misma regla aplicada a dos puntos de referencia distintos.
\end{nota}

\subsection{Derivación del valor del SL}

El SL no se pone en $P_i$ (breakeven): se pone a distancia $B$ bajo el precio del
instante en que se activa.

\begin{align}
b_{\text{efectivo}} &= B \cdot M \\
sl_{\text{nuevo}} &= P_0 - \frac{b_{\text{efectivo}}}{L} = P_0 - \frac{BM}{MmU} = P_0 - \frac{B}{mU}
\end{align}

Evaluando en el instante de activación ($P_0 = P_i + A/(mU)$):

\begin{center}
\colorbox{cfondo}{$\displaystyle SL_1 \;=\; P_i + \frac{A}{mU} - \frac{B}{mU} \;=\; P_i + \frac{A-B}{mU}$}
\end{center}

\begin{nota}
Como $A > B$ en la configuración vigente ($3 > 1{,}5$), $(A-B)/(mU) > 0$ siempre
$\Rightarrow SL_1 > P_i$: el primer SL queda garantizado por encima del precio de
entrada —ya "ganador"— por construcción algebraica, no por casualidad de los valores
actuales.
\end{nota}

\subsection{Ejemplos por activo}

Mismos $P_i$ ilustrativos que la Tabla~1 (AMZN con el valor real de la sesión), para
que se vea directamente que $P_0^{\,SL_1} = P_0^{\,O0\to OE}$ activo por activo —el
punto no es que sea coincidencia, sino que es la misma fórmula.

\begin{table}[h]
\centering
\small
\begin{tabularx}{\linewidth}{@{}l r r r r r r@{}}
\toprule
\textbf{Activo} & $(A-B)/(mU)$ & $B/(mU)$ & $P_i$ & $P_0^{\,SL_1}$ & $SL_1$ & $P_0-SL_1$ \\
\midrule
BTCUSD & 150,00 & 150,00 & 65\,000,00 & 65\,300,00 & \textbf{65\,150,00} & 150,00 \\
ETHUSD & 15,00  & 15,00  & 3\,000,00  & 3\,030,00  & \textbf{3\,015,00}  & 15,00  \\
TSLA   & 1,50   & 1,50   & 250,00     & 253,00     & \textbf{251,50}     & 1,50   \\
GOOGL  & 1,50   & 1,50   & 170,00     & 173,00     & \textbf{171,50}     & 1,50   \\
NVDA   & 1,50   & 1,50   & 130,00     & 133,00     & \textbf{131,50}     & 1,50   \\
AMZN   & 1,50   & 1,50   & 248,37     & 251,37     & \textbf{249,87}     & 1,50   \\
\bottomrule
\end{tabularx}
\caption{Precio de activación del primer SL y valor en que queda. Última columna:
$P_0 - SL_1 = B/(mU)$, la misma distancia que \code{trailing\_stop} mantendrá de ahí
en adelante cada vez que mueva el SL al alza.}
\end{table}

Lectura AMZN (posición real de la sesión, $P_i = 248{,}37$): recién con el precio en
251,37 esa posición recibe su primer SL, y queda puesto en 249,87 — 1,50 por encima de
la entrada y 1,50 por debajo del precio del instante en que se activó.

% ═══════════════════════════════════════════════════════════════════════════
\section{Síntesis}
% ═══════════════════════════════════════════════════════════════════════════

\begin{itemize}[leftmargin=1.4em,itemsep=3pt]
\item Los dos precios de activación comparten la misma forma cerrada:
  $P_i + A/(mU)$. Lo único que cambia es qué es $P_i$: el precio de un soporte
  (Caso~1) o el precio de apertura de una posición (Caso~2).
\item $M$ (\code{LOTAJES\_M}) no aparece en ninguno de los dos umbrales de
  \emph{activación} — se cancela en el Caso~2 y se sacó a propósito del Caso~1. Sí
  sigue afectando \emph{cuánto USD} representa cada evento ($L$ real de la posición),
  vía $B$, \code{PERDIDA\_MAX} y el tamaño de la exposición — pero no el precio al que
  ocurre.
\item El valor del primer SL, $P_i + (A-B)/(mU)$, y la distancia que el trailing
  mantiene después, $B/(mU)$, tampoco dependen de $M$ por la misma razón.
\end{itemize}

\end{document}
